% !TEX TS-program = lualatex

\documentclass[border = 3pt]{standalone}

\usepackage{main}

%%%
% La mise en forme du triangle "isocèle" des combinaiisons est
% aisée via le package \tikz.
%%%
\usepackage{tikz}

\begin{document}

%%%
% caution::
%     La mise en forme reste correcte jusqu'à ''nmax = 12''.
%%%
\def\nmax{10}

\begin{tikzpicture}[
  scale = 0.8,
  cpn-style/.style = {
    rectangle, draw,
    minimum size = 2em,
    fill         = gray!20,
  }
]
  \foreach \n in {0,...,\nmax} {
    \pgfmathsetmacro\pshift{0.5 * \n}

    \foreach \p in {0,...,\n} {
      \draw (\p - \pshift, 6 - \n)
            node[cpn-style] { \Cpn{\p}{\n} };
    }
  }
\end{tikzpicture}

\end{document}
